\documentclass[12pt]{article}
\usepackage[T1]{fontenc}
\usepackage[utf8]{inputenc}
\usepackage{lmodern}
\usepackage{textcomp}
\usepackage{enumitem}
\usepackage{geometry}
\geometry{margin=1in}
\begin{document}
\begin{center}
Answer Key - Sociology - Undergraduate Level Assessment 25 \
\small (Answers are ordered exactly as in the question paper.)
\end{center}

\section*{Multiple Choice}
\begin{enumerate}[label=\arabic*.]
\item \textbf{Answer:} A
\\ \textit{Options:} A) there is a large 'dark figure' of unreported incidents; B) the changing definitions of legal categories have made it harder to convict offenders; C) researchers are not allowed access to official statistics; D) there is no valid or reliable way of researching such a sensitive topic
\\ \textit{Note:} The correct answer is accurate because the term 'dark figure' refers to the significant number of domestic violence incidents that go unreported due to factors such as fear, shame, and lack of support, leading to an incomplete understanding of the true prevalence of the issue.
\item \textbf{Answer:} A
\\ \textit{Options:} A) highly specialized, interrelated sets of social practices; B) disorganized social relations in a postmodern world; C) virtual communities in cyberspace; D) no longer relevant to sociology
\\ \textit{Note:} The correct answer reflects the understanding that in contemporary societies, social institutions, such as family, education, and economy, are characterized by their specialization in specific functions while remaining interconnected and influencing one another through various social practices.
\item \textbf{Answer:} B
\\ \textit{Options:} A) most of the characters in soap operas are women; B) they represent images of women as both domesticated and independent; C) they alienate women and appeal to an audience of men; D) female television producers are most likely to work in this area
\\ \textit{Note:} Pilcher's identification of soap operas as a 'feminine genre' stems from their portrayal of women navigating both traditional roles of domesticity and contemporary ideals of independence, reflecting the complexities of women's identities in society.
\item \textbf{Answer:} B
\\ \textit{Options:} A) race is an objective way of categorizing people on biological grounds; B) the idea of race is socially constructed through powerful ideologies; C) race relations in Britain and America can be traced back to colonial times; D) people choose their racial identity and this becomes fixed
\\ \textit{Note:} The correct answer highlights that theories of racialized discourse emphasize how societal beliefs and power structures shape our understanding of race, indicating that race is not a biological fact but a concept formed through social interactions and ideologies.
\item \textbf{Answer:} D
\\ \textit{Options:} A) the privatization of public services by the Conservative government; B) the lifestyle practice of shopping in peer groups; C) the form of tuberculosis suffered by those who collect stamps; D) the provision of health, housing, and education services by the state
\\ \textit{Note:} The term 'collective consumption' refers to the state's role in providing essential services such as health, housing, and education, highlighting how these services are collectively consumed and organized within a society.
\end{enumerate}

\section*{True / False}
\begin{enumerate}[label=\arabic*.]
\setcounter{enumi}{5}
\item \textbf{Answer:} True
\\ \textit{Options:} A) True; B) False
\\ \textit{Note:} The correspondence principle suggests that the structure and culture of schools mirror the hierarchical and conformist nature of the workforce, thereby socializing students to accept authority and adhere to societal norms essential for their future roles in the labor market.
\item \textbf{Answer:} True
\\ \textit{Options:} A) True; B) False
\\ \textit{Note:} Telework is defined as a work arrangement where employees perform their job duties remotely, often from home, using digital communication tools to connect with their employer and colleagues, which confirms that it is indeed a form of employment involving an outside employer and information technology.
\item \textbf{Answer:} True
\\ \textit{Options:} A) True; B) False
\\ \textit{Note:} A population pyramid with a higher birth rate than death rate typically appears wider at the bottom than at the top because the larger number of younger individuals reflects the higher fertility levels, resulting in a broader base compared to the narrower top representing older age cohorts.
\item \textbf{Answer:} False
\\ \textit{Options:} A) True; B) False
\\ \textit{Note:} The concept of 'power elite,' as described by sociologist C. Wright Mills, refers to a small group of political, corporate, and military leaders who hold significant power in society, but it does not specifically denote the exploitation of the proletariat, which is a characteristic of Marxist theory rather than Mills' framework.
\item \textbf{Answer:} False
\\ \textit{Options:} A) True; B) False
\\ \textit{Note:} Cultural restructuring often involves not only the preservation of traditional landscapes but also the integration of economic transformations that accommodate contemporary needs and practices.
\end{enumerate}

\section*{Long Form Response}
\begin{enumerate}[label=\arabic*.]
\setcounter{enumi}{10}
\item \textbf{Answer:} Sullivan's (2000) study on the dynamics of housework distribution reveals significant patterns influenced by the employment status of men and their partners. The findings indicate that men are more likely to engage in housework when they are unemployed, as their increased availability allows them to take on more domestic responsibilities. Conversely, when both partners work full time, the distribution of housework tends to remain unequal, often with women shouldering a larger share despite their professional commitments. This suggests that traditional gender roles still play a substantial role in domestic labor, highlighting the influence of employment on household dynamics. Overall, Sullivan's research underscores the complex interplay between work status and the allocation of housework, reflecting broader societal norms and expectations.
\\ \textit{Note:} The answer correctly identifies the correlation between men's employment status and their involvement in housework, illustrating that unemployed men tend to take on greater domestic responsibilities, while full-time employment for both partners often perpetuates an unequal distribution of housework, particularly favoring women.
\item \textbf{Answer:} Domhoff identifies several decision-making processes within American society, primarily focusing on the influence of elite groups and the role of institutions in shaping policies and societal norms. He emphasizes the importance of power dynamics and networks among the wealthy, which play a crucial role in determining outcomes in governance and economics. However, Domhoff notably excludes the exploitation process from his framework, which is critical in understanding how systemic inequalities are perpetuated through labor relations and economic structures. This exclusion may limit the comprehensiveness of his analysis, as the exploitation process directly impacts the marginalized populations that are often the subject of elite decision-making. By neglecting this aspect, Domhoff's framework may overlook the foundational economic disparities that inform social stratification and power relations in American society.
\\ \textit{Note:} correct because it accurately captures Domhoff's focus on elite decision-making processes and power dynamics in American society while also recognizing the significant oversight of the exploitation process, which is essential for understanding systemic inequalities.
\end{enumerate}

\vfill
\noindent\textit{End of Answer Key}
\end{document}